
% --- SLIDE 3: TESTING MODELS ---
\begin{frame}{Testing the Model}

    \centering
    \renewcommand{\arraystretch}{1.3}
    {\scriptsize
    \begin{tabular}{p{3.0cm}|p{8.4cm}}
        \toprule
        \textbf{Test type} & \textbf{What it asserts} \\
        \midrule
        \textbf{Smoke / shape}   & The model loads, accepts one row, returns a probability in $[0,1]$. Catches
                                   most deployment breakage. \\
        \textbf{Minimum functionality} & Cases so easy that failure is unambiguous: a customer who cancelled
                                   last month must score high risk. \\
        \textbf{Invariance}      & Perturbations that must \emph{not} change the prediction: reordering
                                   columns, changing a customer ID, whitespace in a category. \\
        \textbf{Directional}     & Perturbations with a known sign: raising the price should not \emph{lower}
                                   predicted churn. \\
        \textbf{Golden set}      & A fixed, curated set of examples with expected outputs, reviewed by a human
                                   and version-controlled. \\
        \textbf{Regression / slice} & Overall metric above a floor \emph{and} no important segment worse than
                                   the incumbent by more than $\delta$. \\
        \bottomrule
    \end{tabular}}

    \vspace{0.21em}
    \begin{itemize}
        \scriptsize
        \item \textbf{Slice tests matter more than the headline metric.} An overall $+1\%$ that hides
              $-8\%$ for one region is a regression, not an improvement.
        \item {\scriptsize The invariance/directional/minimum-functionality vocabulary is from Ribeiro et al.,
              ``Beyond Accuracy: Behavioral Testing of NLP Models with CheckList,'' ACL 2020,
              \href{https://arxiv.org/abs/2005.04118v1}{arXiv:2005.04118v1} --- the idea transfers directly to
              tabular models.}
    \end{itemize}

\end{frame}
